%!TEX TS-program = lualatex
%!TEX encoding = UTF-8 Unicode

\documentclass[12pt, addpoints]{exam}

%\printanswers

\usepackage{fontspec}
\setmainfont[Ligatures={Common, TeX}, BoldFont={* Bold}, ItalicFont={* Italic}, Numbers={Proportional}]{Linux Libertine O}
\setsansfont[Scale=MatchLowercase,Ligatures=TeX]{Linux Biolinum O}
\setmonofont[Scale=MatchLowercase]{Inconsolatazi4}
\usepackage{microtype}


\usepackage{graphicx}
	\graphicspath{{/Users/goby/pictures/teach/438/exam/}} % set of paths to search for images

\usepackage[left=0.75in,right=0.75in,top=1in]{geometry}    

\usepackage{multicol}
\usepackage{amsmath}

\def\short#1{\csname @gobbletwo\expandafter\endcsname\number#1}

\pagestyle{headandfoot}
%\firstpageheader{\bfseries{Ichthyology, Exam 1, F14}}{}{\bfseries{Name: \enspace \makebox[2.5in]{\hrulefill}}}
\firstpageheader{Biogeography, Exam 2, F\short{\year}}{}{%
	\ifprintanswers\textbf{KEY\ \numpoints\ pts}\else Name: \enspace \makebox[2.5in]{\hrulefill}\fi}
\runningheader{}{}{\small{pg. \thepage}}
%\runningheadrule

\footer{}{}{}%{\makebox[0.75in]{\hrulefill}}
\extrafootheight{-0.5in}

\pointsinmargin
\marginpointname{ pts}


%% Remove comment %% to print answer key.
\unframedsolutions
\renewcommand{\solutiontitle}{}
\SolutionEmphasis{\bfseries}

\newcommand*\AnswerBox[2]{%
	\parbox[t][#1]{0.92\textwidth}{%
		\begin{solution}#2\end{solution}}
	\vspace*{\stretch{1}}
}

\newenvironment{AnswerPage}[1]
{\begin{minipage}[t][#1]{0.92\textwidth}%
		\begin{solution}}
		{\end{solution}\end{minipage}
	\vspace*{\stretch{1}}}

\newlength{\basespace}
\setlength{\basespace}{5\baselineskip}

\begin{document}

\begin{questions}

\question[10]
On the figure below, identify the eight terrestrial biogeographic regions. Draw horizontal lines that show the approximate north/south boundaries between each of the regions. Do not draw east/west boundaries. Identify the location of Wallace's Line.

\ifprintanswers
	\includegraphics[width=\textwidth]{globalmap_answer}	
\else
	\includegraphics[width=\textwidth]{GlobalMap}
\fi
\question[5]
Describe the albedo effect and how it contributed to expansion of Pleistocene glaciers.

\AnswerBox{2\baselineskip}{Albedo is amount of light reflected away from earth. More glaciers would have reflected back more sun, leading to cooler temperatures.}

\question[5]\label{ques:faunal_similarity}
\begin{multicols}{2}
This figure shows the similarity of coastal bivalves (clams and mussels) between North America and Europe. Briefly explain how and why bivalve similarity between North America and Europe is changing over time. (\textsc{mya} = millions of years ago.)\vspace{2ex}
	
\AnswerBox{1.5\basespace}{\parbox{0.45\textwidth}{
Similarity decreases with distance between the continents because continued speciation over time causes the continental faunas to become increasingly different.}}

\columnbreak

\includegraphics[width=0.4\textwidth]{atlanticbasin}

\end{multicols}
	
\newpage

%	\vspace{2ex}
	
\question[10]
Briefly explain seafloor spreading and how that process contributes to the pattern displayed in the figure shown in question~\ref{ques:faunal_similarity}. You may draw a picture to aid your explanation.

\AnswerBox{\basespace}{
Seafloor spreading is moving the basins apart. Fauna become more different over time.
}


\question[15]
\begin{multicols}{2}
This figure shows the distribution of three species of pika (rabbit relatives) in the Nearctic (stars and circles) and Palearctic (triangles). Write a clear and detailed hypothesis that includes \textit{one or more} refugia to explain this distribution and how each might have contributed to speciation. Include the proper name for each refugium and assume the pikas evolved from a single common ancestor.

\columnbreak

\hfill \includegraphics[width=0.45\textwidth]{pika_distribution}

\end{multicols}

\AnswerBox{\basespace}{Berengian and southern refugia.}

\newpage


\question[15]
Australia was not covered by glaciers but climate change during the Pleistocene still caused range contraction and expansion for many taxa.  An Australian gecko, \textit{Heteronotia binoei}, has the distribution indicated by the black dots in the inset below.  The large gap in the distribution suggests one of two scenarios: post-Pleistocene range expansion or a long-term barrier to dispersal between the two regions. Strasburg et al. (2007) tested these alternate hypotheses. Below is a simplified version of their parsimony network.  Based on the network, state which hypothesis is supported and describe how the network supports the hypothesis you chose. 

Each circle represents one allele. Shading corresponds to the population sampled; grey is the western cluster and white is the central cluster. Each line connecting two circles indicates one mutation difference between alleles.  Small dashes indicate additional mutational differences between alleles.

{\centering
	\includegraphics[width=\textwidth]{gecko_network}\par
}

\AnswerBox{\basespace}{
The western (grey) clade shows considerable genetic diversity, suggesting it is older and was likely a refuge. The central (white) clade has little genetic diversity, suggesting it is much younger and therefore the result of recent range expansion.
}

\newpage

\begin{parts}
	\part[5] 
	
	\begin{multicols}{2}
		What is the \emph{most likely} distribution of the five families of freshwater crabs shown at right? In one sentence, why?
		
		\ifprintanswers
		 \textbf{Gondwanan}
		\fi
		
		\columnbreak
		
		\includegraphics[width=0.49\textwidth]{crabdistrib}
		
	\end{multicols}

	\vspace{2\baselineskip}
	
	\part[10] 
	Klaus et al. (2011) examined whether the evolutionary relationships of the freshwater crabs were consistent with the distribution you just named.  Below left is a dendrogram that shows the order of land mass breakup.  Below right is a phylogeny for four of the crab families. Compare the dendrogram and phylogeny, and then explain whether the relationship of the crabs is consistent with a vicariant distribution. Provide two specific examples from the two trees to justify your answer. Numbers in trees indicate time of divergence (\textsc{mya}). 
	
	\vspace{\baselineskip}
	
	\includegraphics{crabgondwana.png}
\end{parts}

\begin{solution}
	
	Overall, no. Consistent if you allow for Dispersal from Indomalaysia to Eurasia, placing most Indomalaysian taxa as basal to clade of Neotropics and Africa.  However, still have problems with Socotra and Madagascar.  Madagascar taxa are more closely related to African taxa. Madagascar taxa should be more closely related to Indomalaysian taxa. Madagascar taxa diverged from African taxa circa 76–73 MYA, well after the separation of Madagascar (and India) from Africa. Neotropical taxa more closely related to Socotra and Eurasia taxa rather than African taxa. Seychelles more closely related to one clade of African taxa rather than Indomalaysian/Madagascar taxa.
\end{solution}

\end{questions}

\end{document}